\par\Needspace{7\baselineskip}\bigskip
\begin{center}
\pdfbookmark[2]{Gauss to Olbers, Göttingen, 23 April 1814}{pallas-brief-18140423}
\textsc{Gauss} to \textsc{Olbers}, Göttingen, 23 April 1814.
\end{center}
\noindent\hspace*{\parindent}. . . . .\quad I have already written to you about how closely the 9 oppositions of Pallas agree with my perturbation calculations. Still, the differences remain too large for my liking. The perturbations by Saturn and Mars may perhaps reduce them a little further, but from a rough estimate I have made, the reduction will be slight. It does seem somewhat questionable to give weight to results whose determination required numerical calculation of about 1000 equations for each opposition. Yet all 9000 of these calculations were done twice, and I took the greatest care in calculating the equations themselves. I therefore consider it hardly doubtful that my value for Jupiter's mass still needs correction. To find the true value of that correction, it is essential first to remove the perturbations by Saturn and Mars. \textsc{Nicolai} has nearly finished the former, and I have begun the latter. Still, I could not resist my impatience to see how far those differences could be reduced (still without taking the perturbations by {\AstroSymbols ♄} \SourceLeaf{429}and {\AstroSymbols ♂} into account) if Jupiter's mass itself were treated as an unknown quantity. To my great surprise I found that the agreement improves to a \textls[180]{quite extraordinary} degree. Here are the differences in longitude
\begin{center}
\begingroup\small\setlength{\tabcolsep}{3pt}
\begin{tabular}{@{}rrrrrrrrc@{}}
\multicolumn{9}{l}{with the old mass of {\AstroSymbols ♃}}\\[2pt]
\(-0''{,}5\)&\(-21''{,}0\)&\(-13''{,}2\)&\(+19''{,}6\)&\(-25''{,}2\)&\(+16''{,}6\)&\(+4''{,}5\)&\(-27''{,}7\)&—\\[4pt]
\multicolumn{9}{l}{with the corrected mass}\\[2pt]
\(+3{,}5\)&\(-3{,}1\)&\(+2{,}3\)&\(-0{,}5\)&\(+1{,}3\)&\(-5{,}9\)&\(+1{,}9\)&\(-7{,}1\)&—
\end{tabular}\endgroup
\end{center}
\begin{center}
\begingroup\small\setlength{\tabcolsep}{2pt}
\begin{tabular}{@{}cccccccr@{}}
mean&error&without&correction&of the&mass&. . . . . .&\(17''{,}3\)\\
»&»&with&»&»&»&. . . . . .&\(4{,}3\).
\end{tabular}\endgroup
\end{center}
\par By mean error I mean the square root of the mean value of the square of the error, i.\,e.\ of the sum of the squares divided by their number. The sum of the squares is
\begin{center}
\begin{tabular}{@{}lcr@{}}
without correction&. . . . . .&2684\\
with correction&. . . . . .&167.
\end{tabular}
\end{center}
\par The latitudes, too, which already agreed quite well without the correction, agree a little better with one another after applying it:
\begin{center}
\begin{tabular}{@{}lcccr@{}}
[without correction]&mean&error&. . . . . .&\(5''{,}5\)\\
{}[with correction]&»&»&. . . . . .&\(4{,}0\).
\end{tabular}
\end{center}
\par The correction to the mass is \(\frac{1}{43}\) of the old value, and is an increase. This result is admittedly still uncertain, but I suspect that, once {\AstroSymbols ♄} and {\AstroSymbols ♂} have been taken into account, it will still be at least \(\frac{1}{60}\). I hear that \textsc{Laplace} himself, in edition IV of the Expos. (which I have yet to see), increased the old mass somewhat on the basis of the observations of Saturn, though only very slightly. I believe that Pallas (especially if more oppositions can first be used) must yield more reliable results here. I should like the above communication to remain between us for the present. The old determination was based solely on the 4th satellite of Jupiter, whose distance \textsc{Laplace} took as \(8'16''\). \textsc{Triesnecker} found \(1\frac{1}{2}''\) more, and one would have to add another \(1\frac{1}{2}''\) to obtain my value for the mass.
\par
